% ===========================================================================
%  Jàng Afrig — Lab 1 : Ma première page HTML
%  Construire la page d'accueil en HTML brut
%  Licence 1 DUT1-Informatique — ESP/UCAD — 2025–2026
%  Dr. El Hadji Bassirou TOURÉ
%  Compilé avec : xelatex (2–3 passes)
% ===========================================================================
\documentclass[11pt,a4paper]{article}

\usepackage{fontspec}
\setmainfont{Latin Modern Roman}
\setsansfont{Latin Modern Sans}
\setmonofont{Latin Modern Mono}
\usepackage{polyglossia}
\setdefaultlanguage{french}

\usepackage[margin=2.5cm]{geometry}
\usepackage{amsmath,amssymb}
\usepackage{graphicx}
\graphicspath{{figures/}}
\usepackage{xcolor}
\usepackage{booktabs}
\usepackage{hyperref}
\usepackage{tcolorbox}
\tcbuselibrary{theorems,breakable,skins}
\usepackage{enumitem}
\usepackage{fancyhdr}
\usepackage{array}
\usepackage{pifont}
\usepackage{longtable}
\usepackage{multirow}

\newcommand{\cmark}{\ding{51}}
\newcommand{\xmark}{\ding{55}}

\hypersetup{
    colorlinks=true,
    linkcolor=blue!60!black,
    urlcolor=blue!60!black,
    citecolor=blue!60!black
}

\newtcolorbox{objectifbox}[1][]{colback=blue!5,colframe=blue!50,fonttitle=\bfseries,title=#1,breakable}
\newtcolorbox{infobox}[1][]{colback=green!5,colframe=green!50!black,fonttitle=\bfseries,title=#1,breakable}
\newtcolorbox{attentionbox}[1][]{colback=orange!5,colframe=orange!70,fonttitle=\bfseries,title=#1,breakable}
\newtcolorbox{retenirbox}[1][]{colback=yellow!10,colframe=yellow!50!black,fonttitle=\bfseries,title=#1,breakable}
\newtcolorbox{prereqbox}[1][]{colback=cyan!5,colframe=cyan!50!black,fonttitle=\bfseries,title=#1,breakable}
\newtcolorbox{projetbox}[1][]{colback=teal!5,colframe=teal!50!black,fonttitle=\bfseries,title=#1,breakable}

% --- Code source HTML5 ---
\usepackage{listings}
\lstdefinelanguage{HTML5}{
  morekeywords={<!DOCTYPE,html,head,meta,title,body,h1,h2,h3,h4,p,ul,ol,li,table,tr,th,td,a,img,strong,em,br,div,span},
  sensitive=true,
  morecomment=[s]{<!--}{-->},
  morestring=[b]",
}
\lstset{
  basicstyle=\ttfamily\small,
  backgroundcolor=\color{gray!8},
  frame=single,
  framerule=0.5pt,
  rulecolor=\color{gray!40},
  breaklines=true,
  breakatwhitespace=true,
  showstringspaces=false,
  tabsize=2,
  xleftmargin=4pt,
  xrightmargin=4pt,
  aboveskip=8pt,
  belowskip=8pt,
  literate={é}{{\'e}}1 {è}{{\`e}}1 {ê}{{\^e}}1 {à}{{\`a}}1 {ù}{{\`u}}1
           {ô}{{\^o}}1 {î}{{\^i}}1 {ç}{{\c{c}}}1 {É}{{\'E}}1
           {→}{{$\rightarrow$}}1
}

\pagestyle{fancy}
\fancyhf{}
\fancyhead[L]{\small Jàng Afrig --- Lab 1 : Ma première page HTML}
\fancyhead[R]{\small UCAD --- L1 DUT1-Info}
\fancyfoot[C]{\thepage}
\fancyfoot[L]{\small Dr. El Hadji Bassirou TOURÉ}
\fancyfoot[R]{\small 2025--2026}
\renewcommand{\headrulewidth}{0.4pt}
\renewcommand{\footrulewidth}{0.4pt}

\title{%
\vspace{-1cm}
\textbf{\LARGE Fondamentaux du Web}\\
\textbf{\LARGE HTML, CSS, JavaScript}\\[0.8cm]
\textbf{\huge Lab 1 --- Ma première page HTML}\\[0.3cm]
\textbf{\Large La page d'accueil de Jàng Afrig en HTML brut}\\[0.5cm]
\large Licence 1 DUT1-Informatique --- 2025--2026}
\author{Dr. El Hadji Bassirou TOURÉ\\
Département Génie Informatique\\
Université Cheikh Anta Diop de Dakar}
\date{}

\begin{document}

\maketitle
\thispagestyle{empty}

\vspace{0.5cm}

% ===========================================================================
% BOÎTE PRÉREQUIS
% ===========================================================================
\begin{prereqbox}[Informations pratiques]
\begin{center}
\begin{tabular}{ll}
\toprule
\textbf{Type} & Lab guidé \\
\textbf{Séance} & S2 --- premier Lab du semestre \\
\textbf{Prérequis} & Lab 0 (environnement OK) + CM1 (notions HTML) \\
\textbf{Validation} & Page d'accueil affichée avec Live Server + tag \texttt{v0.1} sur GitHub \\
\textbf{Rapport} & 1-2 pages : ce que j'ai fait + difficultés + lien repo \\
\bottomrule
\end{tabular}
\end{center}
\end{prereqbox}

\begin{objectifbox}[Objectifs du Lab]
À la fin de ce Lab, vous aurez :
\begin{enumerate}[nosep]
    \item Créé le fichier \texttt{index.html} dans votre repo \texttt{jang-afrig}
    \item Écrit la \textbf{structure HTML complète} de la page d'accueil
    \item Utilisé les balises \texttt{h1}, \texttt{h2}, \texttt{p}, \texttt{strong}, \texttt{em}
    \item Construit une \textbf{liste imbriquée} pour présenter les modes de jeu
    \item Construit un \textbf{tableau} (palmarès) avec \texttt{th} et \texttt{td}
    \item Ajouté un \textbf{lien} et une \textbf{image} (drapeau africain)
    \item Posé le tag \textbf{\texttt{v0.1}} sur GitHub
    \item Rédigé votre rapport de Lab
\end{enumerate}
\end{objectifbox}

\tableofcontents

\newpage

% ===========================================================================
\section{Avant de coder}
% ===========================================================================

\subsection{Rappel : le CM1 et le Lab 0}

Au CM1, vous avez vu :
\begin{itemize}[nosep]
    \item Une page web = \textbf{HTML} (structure) + \textbf{CSS} (déco) + \textbf{JS} (comportement)
    \item Une page HTML a un \texttt{<head>} (invisible) et un \texttt{<body>} (visible)
    \item Les balises essentielles : titres (\texttt{h1}-\texttt{h6}), paragraphes, listes, tableaux, liens, images
\end{itemize}

Au Lab 0, vous avez :
\begin{itemize}[nosep]
    \item Installé VS Code + Live Server + Git
    \item Créé votre repo GitHub \texttt{jang-afrig}
    \item Cloné ce repo sur votre ordinateur
    \item Fait un premier cycle Git complet
\end{itemize}

\textbf{Aujourd'hui}, on combine les deux : on écrit une vraie page HTML
dans le repo, on la voit dans le navigateur, on la pousse sur GitHub avec
un tag de version.

\subsection{Ce qu'on va construire concrètement}

La \textbf{page d'accueil} de Jàng Afrig en HTML brut. Pas de CSS (la page
sera moche, c'est normal --- on s'occupera du style aux Labs 3 et 4). Pas
de JavaScript (la page ne fera rien d'interactif --- on s'en occupe à
partir du Lab 5).

Voici ce que la page contiendra :

\begin{itemize}[nosep]
    \item Un titre principal \textbf{Jàng Afrig} avec un sous-titre
    \item Un paragraphe de présentation
    \item La \textbf{liste des modes de jeu} (capitales, drapeaux, langues, voisins, chrono) avec des sous-éléments
    \item Un \textbf{tableau} «\,Palmarès des champions\,» (5 lignes fictives)
    \item Une image (drapeau de l'Union Africaine, ou un drapeau de votre choix)
    \item Un lien vers le repo GitHub du projet
\end{itemize}

\begin{retenirbox}[Pas de panique si c'est moche]
À la fin du Lab, votre page ressemblera à un document Word de 1990 :
texte noir sur fond blanc, sans mise en forme. \textbf{C'est normal et
c'est voulu}. On apprend à \emph{structurer} avant d'apprendre à
\emph{décorer}. La beauté arrive aux Labs 3 et 4 (CSS).
\end{retenirbox}

\newpage

% ===========================================================================
\section{Préparer son espace de travail}
% ===========================================================================

\subsection{Ouvrir VS Code dans le repo}

Dans votre terminal, allez dans votre dossier \texttt{cours-web/jang-afrig}
et ouvrez VS Code :

\begin{lstlisting}
cd ~/cours-web/jang-afrig
code .
\end{lstlisting}

VS Code s'ouvre avec votre repo dans l'Explorateur à gauche. Vous devez
voir au moins :
\begin{itemize}[nosep]
    \item \texttt{README.md} (créé automatiquement par GitHub au Lab 0)
    \item \texttt{index.html} (créé par vous au Lab 0 avec le «\,Hello\,»)
\end{itemize}

\subsection{Faire le ménage dans index.html}

\textbf{Ouvrez \texttt{index.html}} (clic dans l'Explorateur). Vous voyez
le code «\,Hello\,» que vous aviez écrit au Lab 0.

\textbf{Effacez tout le contenu} pour repartir d'une feuille blanche. On
va le réécrire pas à pas.

\begin{infobox}[Pourquoi tout effacer ?]
On aurait pu modifier le fichier petit à petit, mais pour bien intégrer la
structure HTML, c'est mieux de la réécrire \textbf{de zéro}. À la
deuxième page que vous écrirez dans votre vie, vous saurez la taper sans
réfléchir.
\end{infobox}

\subsection{Lancer Live Server}

\textbf{Avant de coder}, lancez Live Server pour voir le résultat en
direct :
\begin{enumerate}[nosep]
    \item Clic droit sur \texttt{index.html} dans l'Explorateur
    \item \textbf{Open with Live Server}
\end{enumerate}

Votre navigateur s'ouvre sur une page \textbf{vide} (puisqu'on a tout
effacé). Laissez le navigateur ouvert à côté de VS Code.

\begin{retenirbox}[Le rituel du développeur]
Pendant tout le Lab, vous travaillerez avec :
\begin{itemize}[nosep]
    \item \textbf{VS Code} à gauche de l'écran
    \item \textbf{Le navigateur} à droite
\end{itemize}
À chaque fois que vous enregistrez (\texttt{Ctrl+S}), la page se rafraîchit
toute seule. Vous voyez le résultat instantanément.
\end{retenirbox}

\newpage

% ===========================================================================
\section{Étape 1 : Le squelette HTML}
% ===========================================================================

On commence par poser le squelette de toute page HTML. Tapez (ou copiez)
ce code dans \texttt{index.html} :

\begin{lstlisting}[language=HTML5]
<!DOCTYPE html>
<html lang="fr">
    <head>
        <meta charset="UTF-8">
        <title>Jàng Afrig</title>
    </head>
    <body>
    </body>
</html>
\end{lstlisting}

Enregistrez (\texttt{Ctrl+S}). Regardez le navigateur : il est toujours
vide \textbf{visuellement}, mais l'\textbf{onglet} affiche désormais
«\,Jàng Afrig\,». \cmark

\begin{infobox}[Pourquoi vide visuellement ?]
Tout ce que vous écrivez dans \texttt{<head>} est invisible. Le \texttt{<body>}
est encore vide. Donc rien à afficher. C'est attendu.
\end{infobox}

\textbf{Ce qui se passe ligne par ligne} :

\begin{center}
\begin{tabular}{@{}lp{8.5cm}@{}}
\toprule
\textbf{Ligne} & \textbf{Rôle} \\
\midrule
\texttt{<!DOCTYPE html>} & «\,Je suis du HTML5 moderne\,» \\
\texttt{<html lang="fr">} & Démarre la page, langue française \\
\texttt{<meta charset="UTF-8">} & Permet les accents (à, é, ç) et le \textbf{Jàng} \\
\texttt{<title>Jàng Afrig</title>} & Texte de l'onglet du navigateur \\
\texttt{<body></body>} & La zone visible, vide pour l'instant \\
\bottomrule
\end{tabular}
\end{center}

\begin{attentionbox}[Le caractère «\,à\,» de Jàng]
Sans \texttt{<meta charset="UTF-8">}, le mot «\,Jàng\,» risque de s'afficher
comme «\,Jng\,» ou «\,J\textbackslash xc3\textbackslash xa0ng\,». Cette ligne est
\textbf{obligatoire} dès qu'on a des accents (donc tout le temps en
français).
\end{attentionbox}

% ===========================================================================
\section{Étape 2 : Titre et présentation}
% ===========================================================================

On va maintenant remplir le \texttt{<body>}. Ajoutez \textbf{entre}
\texttt{<body>} et \texttt{</body>} :

\begin{lstlisting}[language=HTML5]
<h1>Jàng Afrig</h1>
<h2>Apprendre l'Afrique, un pays à la fois</h2>

<p>
    Bienvenue sur <strong>Jàng Afrig</strong>, un quiz géographique sur les
    <strong>54 pays africains</strong>. Choisissez un mode de jeu, répondez à
    10 questions, et battez votre meilleur score !
</p>

<p>
    Le mot <em>jàng</em> signifie <em>«&nbsp;apprendre, étudier&nbsp;»</em>
    en wolof. Le projet est développé dans le cadre du cours
    <strong>Fondamentaux du Web</strong> en L1 DUT1-Info à l'ESP/UCAD.
</p>
\end{lstlisting}

\textbf{Enregistrez}. Le navigateur affiche maintenant un grand titre
«\,Jàng Afrig\,», un sous-titre «\,Apprendre l'Afrique\ldots\,», et deux
paragraphes. \cmark

\begin{infobox}[L'entité \texttt{\&nbsp;}]
\texttt{\&nbsp;} est un \textbf{espace insécable} (qui ne se coupe pas en
fin de ligne). On l'utilise avant les guillemets français
(«\,texte\,»), les deux-points, les points-virgules. Pas obligatoire,
mais joli.
\end{infobox}

\newpage

% ===========================================================================
\section{Étape 3 : La liste des modes de jeu}
% ===========================================================================

C'est ici qu'on va utiliser une \textbf{liste imbriquée} : une liste
principale des modes, et pour chaque mode, une sous-liste avec une
description courte.

\textbf{Ajoutez à la suite}, dans le \texttt{<body>} :

\begin{lstlisting}[language=HTML5]
<h2>Les modes de jeu</h2>

<ul>
    <li>
        <strong>Capitales</strong>
        <ul>
            <li>Identifier la capitale d'un pays africain</li>
            <li>Exemple : « Quelle est la capitale du Mali ? »</li>
        </ul>
    </li>
    <li>
        <strong>Drapeaux</strong>
        <ul>
            <li>Reconnaître le drapeau d'un pays</li>
            <li>Exemple : « À quel pays appartient ce drapeau ? »</li>
        </ul>
    </li>
    <li>
        <strong>Langues</strong>
        <ul>
            <li>Associer une langue à un pays</li>
            <li>Exemple : « Dans quel pays parle-t-on le swahili ? »</li>
        </ul>
    </li>
    <li>
        <strong>Voisins</strong>
        <ul>
            <li>Compter les pays frontaliers</li>
            <li>Exemple : « Combien de voisins a l'Algérie ? »</li>
        </ul>
    </li>
    <li>
        <strong>Chrono</strong>
        <ul>
            <li>Maximum de bonnes réponses en 60 secondes</li>
        </ul>
    </li>
</ul>
\end{lstlisting}

\textbf{Enregistrez}. Vous voyez maintenant la liste à puces avec, sous
chaque mode, une sous-liste indentée. \cmark

\begin{retenirbox}[Indentation : votre meilleur ami]
Notez comme le code est \textbf{indenté} : chaque niveau de liste imbriquée
est décalé de 4 espaces. \textbf{Pour l'ordinateur}, ça ne change rien.
\textbf{Pour vous} (et pour quiconque relit votre code), c'est
\textbf{indispensable} pour comprendre la structure.

VS Code indente automatiquement quand vous appuyez sur \texttt{Tab}.
\end{retenirbox}

\newpage

% ===========================================================================
\section{Étape 4 : Le tableau des champions}
% ===========================================================================

On va maintenant ajouter un \textbf{tableau} avec un faux palmarès. Les
vrais scores arriveront au Lab 10 (avec \texttt{localStorage}), pour
l'instant on hardcode des données.

\textbf{Ajoutez à la suite} :

\begin{lstlisting}[language=HTML5]
<h2>Palmarès des champions</h2>

<p>Les meilleurs scores enregistrés (données fictives pour le moment) :</p>

<table>
    <tr>
        <th>Rang</th>
        <th>Pseudo</th>
        <th>Mode</th>
        <th>Score</th>
        <th>Date</th>
    </tr>
    <tr>
        <td>1</td>
        <td>Fatou</td>
        <td>Capitales</td>
        <td>10/10</td>
        <td>15 mars 2026</td>
    </tr>
    <tr>
        <td>2</td>
        <td>Moussa</td>
        <td>Drapeaux</td>
        <td>9/10</td>
        <td>12 mars 2026</td>
    </tr>
    <tr>
        <td>3</td>
        <td>Aminata</td>
        <td>Capitales</td>
        <td>9/10</td>
        <td>10 mars 2026</td>
    </tr>
    <tr>
        <td>4</td>
        <td>Cheikh</td>
        <td>Voisins</td>
        <td>8/10</td>
        <td>8 mars 2026</td>
    </tr>
    <tr>
        <td>5</td>
        <td>Awa</td>
        <td>Langues</td>
        <td>8/10</td>
        <td>5 mars 2026</td>
    </tr>
</table>
\end{lstlisting}

\textbf{Enregistrez}. Le tableau apparaît, sans bordures (le navigateur
les cache par défaut --- on les verra au Lab 3 avec CSS). \cmark

\begin{infobox}[Tableau sans bordures ?]
Sans CSS, un tableau HTML s'affiche \textbf{sans aucune bordure}. C'est le
comportement standard depuis HTML5. Au Lab 3, deux lignes de CSS suffiront
à dessiner de jolies bordures et alterner les couleurs des lignes.
\end{infobox}

\newpage

% ===========================================================================
\section{Étape 5 : Image et lien}
% ===========================================================================

\subsection{Préparer le dossier images}

Dans votre repo, on va mettre les images dans un sous-dossier
\texttt{images/} pour ranger.

\begin{enumerate}[nosep]
    \item Dans l'Explorateur de VS Code, clic droit sur le dossier racine
    \item \textbf{New Folder}
    \item Nommez-le : \texttt{images}
\end{enumerate}

\subsection{Récupérer une image de drapeau}

Choisissez un drapeau parmi ceux-ci :
\begin{itemize}[nosep]
    \item Drapeau du Sénégal : \url{https://flagcdn.com/w320/sn.png}
    \item Drapeau de l'Union Africaine : \url{https://upload.wikimedia.org/wikipedia/commons/thumb/8/8c/Flag_of_the_African_Union.svg/320px-Flag_of_the_African_Union.svg.png}
    \item Drapeau du Mali : \url{https://flagcdn.com/w320/ml.png}
    \item Drapeau du Maroc : \url{https://flagcdn.com/w320/ma.png}
    \item Drapeau de la Côte d'Ivoire : \url{https://flagcdn.com/w320/ci.png}
\end{itemize}

\textbf{Téléchargez} le drapeau de votre choix : clic droit sur le lien
$\to$ \textbf{Enregistrer la cible sous} $\to$ choisissez le dossier
\texttt{images/} de votre repo, nommez le fichier \texttt{drapeau.png}.

\subsection{Insérer l'image et le lien}

\textbf{Ajoutez à la suite} dans \texttt{<body>} :

\begin{lstlisting}[language=HTML5]
<h2>À propos de Jàng Afrig</h2>

<p>
    <img src="images/drapeau.png" alt="Drapeau africain" width="200">
</p>

<p>
    Le code source du projet est disponible sur GitHub :
    <a href="https://github.com/votre-username/jang-afrig">
        github.com/votre-username/jang-afrig
    </a>
</p>

<p>
    Données fournies par
    <a href="https://restcountries.com">REST Countries</a>,
    une API libre et gratuite sur les pays du monde.
</p>
\end{lstlisting}

\textbf{Remplacez \texttt{votre-username}} par votre vrai \emph{username}
GitHub. Enregistrez. \cmark

\begin{attentionbox}[L'attribut \texttt{alt} n'est pas optionnel]
Toujours mettre un \texttt{alt} sur chaque \texttt{<img>}. Pour vos
camarades qui utilisent un lecteur d'écran (cécité, déficience visuelle),
c'est leur seule façon de «\,voir\,» l'image. C'est aussi utile si
l'image ne se charge pas (lien cassé, connexion lente).
\end{attentionbox}

\newpage

% ===========================================================================
\section{Vérification : ce que doit afficher la page}
% ===========================================================================

À ce stade, votre page doit ressembler à ce schéma (en plus brut sans
style) :

\begin{center}
\fbox{
\begin{minipage}{0.85\textwidth}
\vspace{0.3cm}
\begin{center}
{\Large \textbf{Jàng Afrig}}\\[2pt]
{\large Apprendre l'Afrique, un pays à la fois}
\end{center}

\medskip

\small Bienvenue sur \textbf{Jàng Afrig}, un quiz géographique sur les
\textbf{54 pays africains}\ldots

\medskip

{\large \textbf{Les modes de jeu}}

\begin{itemize}[leftmargin=*,nosep,label=\textbullet]
    \item \textbf{Capitales}
        \begin{itemize}[nosep,label=\textopenbullet]
            \item Identifier la capitale d'un pays africain
            \item Exemple : «\,Quelle est la capitale du Mali ?\,»
        \end{itemize}
    \item \textbf{Drapeaux}\ldots
    \item \textbf{Langues}\ldots
\end{itemize}

\medskip

{\large \textbf{Palmarès des champions}}

\medskip

\begin{tabular}{lllll}
Rang & Pseudo & Mode & Score & Date \\
1 & Fatou & Capitales & 10/10 & 15 mars 2026 \\
2 & Moussa & Drapeaux & 9/10 & 12 mars 2026 \\
\multicolumn{5}{l}{\ldots}
\end{tabular}

\medskip

{\large \textbf{À propos de Jàng Afrig}}

\medskip

\fbox{[Image du drapeau choisi]}

\medskip

\small Le code source du projet est disponible sur GitHub : \\
\textcolor{blue}{\underline{github.com/votre-username/jang-afrig}}\\
\vspace{0.3cm}
\end{minipage}
}
\end{center}

\bigskip

\begin{retenirbox}[Aspect attendu]
Tout est noir, sur fond blanc, sans bordure de tableau, sans couleurs ni
mise en page. \textbf{C'est volontaire}. La page est \emph{structurée}, pas
\emph{décorée}. CSS arrive bientôt.
\end{retenirbox}

\subsection{Si quelque chose ne s'affiche pas comme attendu}

\begin{itemize}[leftmargin=*,nosep]
    \item \textbf{Page totalement blanche} : vérifiez que vous avez bien enregistré (\texttt{Ctrl+S}) et que Live Server tourne.
    \item \textbf{Caractères bizarres} (Jng au lieu de Jàng) : vérifiez la ligne \texttt{<meta charset="UTF-8">}.
    \item \textbf{L'image ne s'affiche pas} : vérifiez le chemin (\texttt{images/drapeau.png}) et que le fichier est bien dans le bon dossier.
    \item \textbf{La sous-liste n'apparaît pas} : vérifiez que le \texttt{<ul>} imbriqué est bien \emph{à l'intérieur} d'un \texttt{<li>}, pas entre deux \texttt{<li>}.
\end{itemize}

\newpage

% ===========================================================================
\section{Exercices d'application}
% ===========================================================================

\begin{objectifbox}[Trois mini-exercices à faire dans l'ordre]
\textbf{Exercice 1 --- Une 6\textsuperscript{e} ligne au tableau}

Ajoutez une ligne supplémentaire au tableau du palmarès, avec :
\begin{itemize}[nosep]
    \item Rang : 6
    \item Pseudo : votre prénom
    \item Mode : un mode de votre choix
    \item Score : un score de votre choix
    \item Date : aujourd'hui
\end{itemize}

\medskip

\textbf{Exercice 2 --- Un mode de jeu supplémentaire}

Inventez un \textbf{6\textsuperscript{e} mode de jeu} et ajoutez-le à la
liste des modes (avec sous-liste de 1 ou 2 lignes pour le décrire). Idées :
«\,Monnaies\,», «\,Superficie\,», «\,Population\,»\ldots

\medskip

\textbf{Exercice 3 --- Une section «\,Contact\,»}

Ajoutez en bas de la page une section :
\begin{itemize}[nosep]
    \item Un titre \texttt{<h2>Contact</h2>}
    \item Un paragraphe avec votre nom complet, votre filière (DUT1-Info), votre email étudiant
    \item Un \emph{lien mailto} vers votre email : \texttt{<a href="mailto:vous@esp.sn">vous@esp.sn</a>}
\end{itemize}

\medskip

\textbf{Bonus} (si vous avez fini en avance) : ajoutez 2 ou 3 drapeaux
supplémentaires dans la section «\,À propos\,», côte à côte. Vous verrez
qu'ils s'affichent les uns en dessous des autres. C'est normal, on
règlera ça avec Flexbox au Lab 4.
\end{objectifbox}

\newpage

% ===========================================================================
\section{Aide-mémoire}
% ===========================================================================

\subsection{Les balises vues dans ce Lab}

\begin{center}
\begin{tabular}{@{}lp{8cm}@{}}
\toprule
\textbf{Balise} & \textbf{Rôle} \\
\midrule
\texttt{<!DOCTYPE html>}    & Déclaration HTML5 (toujours en première ligne) \\
\texttt{<html lang="fr">}   & Racine de la page, langue \\
\texttt{<head>}             & Métadonnées (invisibles) \\
\texttt{<meta charset="UTF-8">} & Encodage des caractères (accents) \\
\texttt{<title>}            & Texte de l'onglet \\
\texttt{<body>}             & Contenu visible \\
\texttt{<h1>} \ldots \texttt{<h6>} & Titres (un seul \texttt{<h1>} par page) \\
\texttt{<p>}                & Paragraphe \\
\texttt{<strong>}           & Texte important (gras par défaut) \\
\texttt{<em>}               & Mise en valeur (italique par défaut) \\
\texttt{<ul>} / \texttt{<ol>} / \texttt{<li>} & Listes non ordonnée / ordonnée / élément \\
\texttt{<table>} / \texttt{<tr>} / \texttt{<th>} / \texttt{<td>} & Tableau / ligne / cellule en-tête / cellule \\
\texttt{<a href="...">}     & Lien hypertexte \\
\texttt{<img src="..." alt="...">} & Image (auto-fermante) \\
\bottomrule
\end{tabular}
\end{center}

\subsection{Le squelette à connaître par cœur}

\begin{lstlisting}[language=HTML5]
<!DOCTYPE html>
<html lang="fr">
    <head>
        <meta charset="UTF-8">
        <title>Titre de la page</title>
    </head>
    <body>
        <!-- contenu visible ici -->
    </body>
</html>
\end{lstlisting}

\begin{retenirbox}[Le réflexe à acquérir]
Quand vous créez un nouveau fichier HTML dans VS Code, tapez
\texttt{html:5} puis \texttt{Tab}. VS Code génère automatiquement le
squelette pour vous (c'est l'autocomplétion Emmet, intégrée).
\end{retenirbox}

\newpage

% ===========================================================================
\section{Routine Git de fin de Lab}
% ===========================================================================

C'est le moment de \textbf{sauvegarder votre travail} et de marquer cette
version. Vous l'avez fait au Lab 0, on refait pareil.

\subsection{Les commandes à enchaîner}

Dans le terminal intégré de VS Code (\texttt{Ctrl+\textasciigrave}) :

\begin{lstlisting}
git status
git add .
git commit -m "v0.1 : page d'accueil HTML brut"
git push
git tag v0.1
git push --tags
\end{lstlisting}

\subsection{Vérification}

\begin{enumerate}[nosep]
    \item Allez sur \url{https://github.com/votre-username/jang-afrig}
    \item Vous devez voir \texttt{index.html} et le dossier \texttt{images/}
    \item Cliquez sur l'onglet \textbf{Tags} (ou \textbf{Releases}) en haut
    \item Le tag \texttt{v0.1} doit apparaître \cmark
\end{enumerate}

\begin{retenirbox}[Pourquoi un tag ?]
Un \textbf{tag Git} est une étiquette qu'on pose sur une version précise
du projet. Si vous travaillez 6 mois plus tard sur le projet et que vous
voulez retrouver \emph{l'état exact à la fin du Lab 1}, il vous suffit
de faire \texttt{git checkout v0.1}. Vous y êtes.
\end{retenirbox}

% ===========================================================================
\section{Rapport à produire}
% ===========================================================================

\begin{attentionbox}[Rapport du Lab 1 --- obligatoire]
À la fin de ce Lab, produisez un \textbf{rapport court} de
\textbf{1 à 2 pages maximum}.

\medskip

\textbf{Format du fichier :}
\begin{itemize}[nosep]
    \item PDF, Markdown (\texttt{.md}) ou Word
    \item Nom : \texttt{rapport-lab-1-NOM-Prenom.pdf}\\
          (ex. \texttt{rapport-lab-1-DIOP-Fatou.pdf})
    \item Déposer dans le sous-dossier \texttt{rapports/} de votre repo GitHub
\end{itemize}

\medskip

\textbf{Contenu obligatoire (4 sections) :}

\begin{enumerate}[nosep]
    \item \textbf{Ce que j'ai fait} (½ page) : page d'accueil créée, balises utilisées, exercices d'application complétés.
    \item \textbf{Difficultés rencontrées} (¼ page) : ce qui m'a bloqué (image qui ne s'affiche pas, structure de liste imbriquée, encodage\ldots).
    \item \textbf{Comment j'ai résolu ces difficultés} (¼ page) : recherche, aide d'un camarade, essais et erreurs.
    \item \textbf{Lien vers mon repo GitHub avec le tag v0.1} : \\
          ex. \url{https://github.com/fatou-diop/jang-afrig/releases/tag/v0.1}
\end{enumerate}

\medskip

\textbf{Ton attendu :} honnête et direct. Si quelque chose n'a pas marché
du premier coup, dites-le. L'erreur fait partie de l'apprentissage.
\end{attentionbox}

% ===========================================================================
\section{Résumé et transition}
% ===========================================================================

\subsection{Ce que vous avez accompli}

\begin{itemize}[leftmargin=*,nosep]
    \item Créé une vraie page HTML structurée \cmark
    \item Maîtrisé le squelette HTML5 (DOCTYPE, head, body) \cmark
    \item Utilisé titres, paragraphes, mise en forme inline \cmark
    \item Construit une liste imbriquée (modes de jeu) \cmark
    \item Construit un tableau (palmarès) \cmark
    \item Inséré une image et des liens \cmark
    \item Posé le tag \texttt{v0.1} sur GitHub \cmark
\end{itemize}

\subsection{Ce qu'on fera au Lab 2}

\begin{itemize}[leftmargin=*,nosep]
    \item Découvrir les balises \textbf{sémantiques} HTML5 (\texttt{header}, \texttt{nav}, \texttt{main}, \texttt{section}, \texttt{footer})
    \item Construire les \textbf{5 pages} de Jàng Afrig (Accueil, Quiz, Résultat, Mes Scores, À propos)
    \item Ajouter des \textbf{formulaires} (saisie du pseudo, sélection du mode)
    \item Lier les pages entre elles avec des liens internes
    \item Tag \texttt{v0.2}
\end{itemize}

\begin{infobox}[Avant le Lab 2]
Lisez le PDF du Lab 2 que le prof vous a distribué. Aucun code à produire
chez vous. Si vous avez du temps, vous pouvez parcourir la documentation
des formulaires HTML sur MDN :\\
\url{https://developer.mozilla.org/fr/docs/Learn/Forms}
\end{infobox}

\vspace{0.5cm}

\begin{center}
\fbox{
\begin{minipage}{0.85\textwidth}
\centering
\vspace{0.5cm}
\textbf{\Large Lab 1 --- Terminé \cmark}\\[0.3cm]
\normalsize
Vous avez écrit votre première vraie page HTML.\\
La structure est posée. Place à la décoration et à l'interactivité.\\[0.2cm]
\textit{Rendez-vous au Lab 2 !}\\[0.3cm]
\vspace{0.3cm}
\end{minipage}
}
\end{center}

\end{document}
